\documentclass[11pt,a4paper]{article} % Define clase y tamaño de documento

% Configuración de idioma
\usepackage[spanish]{babel} % Soporte para español
\usepackage[utf8]{inputenc} % Codificación UTF-8 (opcional en versiones recientes de LaTeX)
\usepackage[T1]{fontenc} % Codificación de fuente

% Márgenes y layout
\usepackage[rmargin=2.54cm,lmargin=2.54cm,top=2.54cm,bottom=2.54cm,headheight=14pt]{geometry} % Márgenes y altura del encabezado
\usepackage{fancyhdr} % Encabezados y pies de página

% Paquetes matemáticos
\usepackage{amsmath,amsfonts,amssymb} % Símbolos y fuentes matemáticas

% Gráficos y figuras
\usepackage{graphicx} % Inclusión de gráficos avanzados
\DeclareGraphicsExtensions{.pdf,.png,.jpg} % Extensiones gráficas permitidas
\usepackage{tikz} % Dibujos y gráficos vectoriales

%NUEVOS DE ESTE DOC

\usepackage{pgfplots}
\usetikzlibrary{decorations.pathreplacing, positioning}


% Otros paquetes útiles
\usepackage{ragged2e} % Justifica títulos
\usepackage{xparse} % Crear comandos personalizados
\usepackage{stackengine} % Apilamiento de objetos
\usepackage{cancel} % Tachar elementos
\usepackage{color} % Colores
\usepackage{apacite} % Citas estilo APA
\usepackage{hyperref} % Enlaces y referencias
\hypersetup{
    colorlinks=false,  %Ahí está para que los links NO tengan color
    linkcolor=blue,  %Igual si quiere que tengan, acá puede cambiar el color
    filecolor=magenta,      
    urlcolor=cyan,
    pdftitle={Documento},
    pdfpagemode=FullScreen,
}

% Mejorar la tipografía
\usepackage{microtype} % Mejoras tipográficas
\usepackage{csquotes} % Citas adecuadas al idioma

%%%%%%%%%%%%%%%%%%%%%%%%%%%%
% Encabezado y pie de página

\pagestyle{fancy} % Activar estilo fancyhdr
\rhead{Nombre} % Encabezado derecho
\lhead{Microeconomía} % Encabezado izquierdo
\fancyfoot[R]{\thepage} % Pie de página derecho con número de página
\fancyfoot[C]{Institución} % Pie de página centrado con texto

\bibliographystyle{apacite} % Estilo de bibliografía APA

\title{Taller Microeconomía} % Título del documento
\author{Nombre} % Autor
\date{\today} % Fecha

\usetikzlibrary{babel} % Compatibilidad TikZ con babel. Esto ha prevenido muchos errores de vainas de tikz que se petan por babel

\begin{document} % Inicio del documento

\maketitle\thispagestyle{fancy} % Título del documento con estilo fancy

%%%%%%%%%%%%%%%%%%%%%%%%%%%%%%%%

\section{}


\subsection{}

\noindent Calculamos las tasas de oportunidad de producción, es decir, la pendiente de la curva de oferta. 

\begin{itemize}
\item País A: $\frac{100-0}{0-200} = \frac{100}{200} = \frac{\text{1 auto}}{\text{2 computadoras}}$
\end{itemize}

\noindent Así, la tasa es de 0,5 autos por computador y 2 computadoras por auto.


\begin{itemize}
    \item País B: $\frac{0-150}{50-0} = \frac{150}{50} = \frac{\text{3 computadores}}{\text{1 auto}}$
\end{itemize}

\noindent Así, la tasa es 3 computadores por auto y 0,33 autos por computadora.

\subsection{}

\begin{tikzpicture}
    \begin{axis}[
        xlabel={Computadoras},
        ylabel={Autos},
        xmin=0, xmax=200,
        ymin=0, ymax=200,
        xtick={0,50,100,150,200},
        ytick={0,50,100,150,200},
        legend pos=north east,
        grid=major,
        width=\linewidth,
        height=0.6\linewidth,
    ]
    % Data for País A
    \addplot[
        color=blue,  % Color azul para País A
        mark=*,
        ]
        coordinates {
        (0, 50)(150, 0)
        };
    \addlegendentry{País B}

    % Data for País B
    \addplot[
        color=red,  % Color rojo para País B
        mark=*,
        ]
        coordinates {
        (0, 100)(200, 0)
        };
    \addlegendentry{País A}

    \end{axis}
\end{tikzpicture}


\subsection{}

\noindent Ambas fronteras de posibilidades de producción tienen pendientes diferentes, lo cual muestra que los países tienen diferencias en sus tasas relativas, y por ende, uno de los dos bienes les resulta más barato de producir en términos relativos respecto al otro país, debido a que tiene ventajas comparativas en su producción se verían beneficiados de una especialización en alguno de los 2 bienes (en aquel que tengan ventaja) y luego entrar en el comercio entre ambos


\subsection{}

\[
\begin{array}{c|cc}
    & \text{Autos} & \text{Computadoras} \\ \hline
    \text{A} & 0,5 & 2 \\
    \text{B} & 0,33 & 3 \\
\end{array}
\]

\noindent Notemos que sacrificando un computador el país A puede hacer más autos que B. Asimismo, sacrificando un auto, B puede hacer más computadoras que A. Luego, le conviene a A especializarse en autos y a B en computadores, y una vez que ambos se hayan especializado, comerciar entre ellos. En el caso de que se especialicen, el punto de A es (100, 0)\footnote{Es decir, 100 autos y 0 computadores}, y el de B es (0, 150). 

Supongamos un caso hipotético donde los términos de intercambio de ambos países son de 2,5 computadores por un auto. Los países podrían comerciar y podrían alcanzar puntos que estén fuera de su frontera de posibilidades de producción, por lo cual están recibiendo ganancias de la especialización y el comercio. Especializarse le permitiría a ambos países alcanzar puntos que en sus condiciones originales de producción no hubieran podido alcanzar. 


\section{}

\subsection{}

\noindent Cálculo de la pendiente curva oferta:

$$
m = \frac{Q_{2}-Q_{1}}{P_{2}-P_{1}} = \frac{200-100}{20\$-10\$} = 10
$$

\subsection{}

\noindent Para encontrar la ecuación de la curva de oferta, sabemos que $Q_{2} = a + bP_{2}$. Como $b=10$, luego con $Q_{2} = 200$ tenemos que $200 = a + 10(20\$)$, de lo que se sigue que $a = 0$. Así, hemos encontrado todos los componentes de la ecuación. 

\subsection{}

\noindent La ecuación resultante es Q = 10P. (completar)

\subsection{}

\noindent Si $P = 30$, entonces $Q=a+b(30)=300$

\section{}

\subsection{} 

\noindent Tenemos que $Q^{d} = 200 - 5P$ y $Q^{s} = 2P$. Busquemos el equilibrio, es decir, $Q^{d} = Q^{s}$. Así, $200 - 5P = 2P$, de donde se sigue que $P^* = \frac{200}{7} \approx 28,571$. Reemplazando $P^*$ para encontrar $Q^*$, $Q^* = 200 \cdot \frac{200}{7} \approx 57,142$ 

\subsection{} 

\noindent El precio máximo de reserva es el precio para el cual $Q^d=0$. Luego, $0=200-5P_{max}$, de lo que se sigue que $P_{max} = 40$.

El excedente del consumidor (EC) está dado por: $\frac{(P_{max}-P^*)(Q^*)}{2}$. Reemplazando los valores, $EC = \frac{(40-(\frac{200}{7}))(\frac{400}{7})}{2} \approx 326,530612$.

\subsection{}

\noindent El excedente del productor está dado por $\frac{(P^{*}-P_{min})(Q*)}{2}$. Como $P_{min} = 0$, entonces $\frac{(P*)(Q*)}{2} \approx 816,326531$. 

\subsection{}

\noindent Como el excedente total (ET) es la suma del excedente del productor mas el excedente del consumidor, $ET \approx 1142,85653$

\section{}

\subsection{}

\noindent Para encontrar el punto de equilibrio, igualamos $Q^d = Q^s$, así, $300-5P=2P$, de donde se sigue que $P* = \frac{300}{7} \approx 42,8571$. Ahora bien, $Q* = 2P* \approx 85,7142857$

\subsection{}

\noindent El precio de reserva $P_{max}$ es aquel que hace que la cantidad demandada sea 0, así

$P_{max} = \frac{300}{5} = 60$. De esta forma, $EC = \frac{(60-\frac{300}{7})(85,714)}{2} \approx 734,693$ 


\subsection{}

\noindent El nuevo excedente del productor está dado por $\frac{(P^{*}-0)(Q^{*})}{2} = \frac{(42,8571)(85,7142857)}{2} \approx 1836,726$

\subsection{}

\noindent Sumando el nuevo excedente del productor y el nuevo excedente del consumidor, tenemos que $ET \approx 2571,419$ 



%\bibliography{citas}

%%%%%%%%%%%%%%%%%%%%%%%%%%%%%%%
\end{document} %Acá termina
